\section{Methodology}
\label{sec:methodology}

\subsection{From piecewise windows to a global GP fit}

The original HPS GPR studies explored piecewise mass windows, explicit sideband
reconfiguration at each hypothesis, and several alternative background-profile
constructions. Those studies were useful for establishing that a GP could interpolate
the smooth trident spectrum, but they also made clear where the older strategy became
fragile. Window-by-window fits produced visibly less smooth limit curves, amplified the
dependence of the result on local fit boundaries, and complicated any attempt to write
a statistically coherent multi-dataset combination. The present workflow therefore uses
a full-range per-hypothesis GP fit over the histogram range of each dataset and
excludes only a narrow mass-dependent training-exclusion region around the test mass. That choice
yields a background description that is smoother, easier to validate, and naturally
compatible with a combined likelihood in a shared parameter of interest.

\subsection{Observable, scan grid, and blinding}

The input observable is the reconstructed $e^+e^-$ invariant-mass histogram for each
dataset. Throughout this section, $d$ labels the dataset, $m$ denotes the reconstructed
$e^+e^-$ invariant mass, $m_0$ is the test-mass hypothesis, and $\sigma_d(m)$ is the
dataset-specific mass resolution introduced in
Section~\ref{sec:event-selection}. The default scan step in the package is
\begin{equation}
\Delta m = \SI{1}{MeV},
\end{equation}
and the blind window for a mass hypothesis $m_0$ is centered on the detector
resolution,
\begin{equation}
\mathcal{B}_d(m_0) =
\left[m_0 - n_{\mathrm{blind}}\sigma_d(m_0),\,
      m_0 + n_{\mathrm{blind}}\sigma_d(m_0)\right].
\end{equation}
The package exposes the extraction half-width and GP training exclusion as independent
configuration values. Earlier injection and GP-toy validation studies used
$n_{\mathrm{blind}}=1.64$ with the same $1.64\,\sigma_d$ training exclusion. The current
observed-limit and refmatched closure baseline documented later in this note uses a
wider $2.25\,\sigma_d$ blind/extraction window with the GP training exclusion also set
to $2.25\,\sigma_d$. For an ideal Gaussian line shape, the $1.64\,\sigma_d$ interval
contains
\begin{equation}
f_{\mathrm{blind}}^{\mathrm{ideal}}(1.64)
= \mathrm{erf}\!\left(\frac{1.64}{\sqrt{2}}\right)
\simeq 0.899,
\end{equation}
while the $2.25\,\sigma_d$ interval contains
\begin{equation}
f_{\mathrm{blind}}^{\mathrm{ideal}}(2.25)
= \mathrm{erf}\!\left(\frac{2.25}{\sqrt{2}}\right)
\simeq 0.976.
\end{equation}
This point matters conceptually and numerically: the current code does
\emph{not} renormalize the signal template inside either blind window. Instead, it
keeps the full Gaussian normalization and lets only the blinded subset enter the local
likelihood.

The notation $\mathcal{B}_d(m_0)$ denotes the set of histogram bins that lie inside the
blinded interval for dataset $d$ at test mass $m_0$. In the current production
configuration there is no additional gap between the extraction window and the GP
training sample: when the extraction and training-exclusion half-widths are equal, the
training exclusion starts immediately outside $\mathcal{B}_d(m_0)$. A wider exclusion,
for example $1.96\,\sigma_d$ instead of $1.64\,\sigma_d$, reduces signal leakage into
the training region but also increases the interpolation distance and therefore the GP
predictive uncertainty. In the simple background-dominated proxy
$S/\sqrt{B}\propto F(k)/\sqrt{k}$ with $F(k)=\mathrm{erf}(k/\sqrt{2})$, the optimum lies
near $k\simeq1.4$, so the historical $1.64\,\sigma_d$ extraction was close to the
pure-sensitivity optimum. Section~\ref{sec:toys} then motivates the current
$2.25\,\sigma_d$ blind/extraction choice as a signal-recovery improvement after
full-refit pseudoexperiments showed appreciable training-region absorption at narrower
exclusions.

The bin assignment is implemented with the fixed histogram bin centers, not by splitting
partly overlapping bins. Let $c_i=(e_i+e_{i+1})/2$ be the center of bin $i$. The full
bin is included in the blinded likelihood when
\begin{equation}
m_0-n_{\mathrm{blind}}\sigma_d(m_0)
\le c_i \le
m_0+n_{\mathrm{blind}}\sigma_d(m_0),
\end{equation}
and the full bin is included in the GP sideband training sample only when its center
lies strictly outside the training-exclusion interval. Thus a bin whose edge overlaps
the blind boundary is treated by its center: if the center is inside the blind interval,
the entire bin is blinded and enters the local likelihood; if the center is outside,
the bin remains available for training unless it is removed by a wider
training-exclusion interval. No fractional observed count is assigned to both regions.
This same center mask is used consistently for the observed counts, the GP prediction
points, the covariance slice, and the blinded-bin signal-template slice.

The finite binning is therefore part of the definition of the numerical result. The
Gaussian signal template is integrated over the full bin edges in
Eq.~\eqref{eq:window-template}, but the selected set of bins changes discretely when a
boundary passes a bin center. Consequently the quoted local significance and
$\eps^2$ limit are not meaningful to a precision finer than the mass binning and scan
grid used to define these masks. The production toys and validation plots use the same
binning convention, so this edge-bin effect is accounted for by applying the identical
discrete procedure to data, pseudoexperiments, background predictions, signal
templates, and yield-to-coupling conversions.

\subsection{Log-space preprocessing and the choice \texorpdfstring{$\alpha=1/y$}{alpha=1/y}}

Following recent HEP GPR background-modeling studies
\cite{BarrLiu2025,Gandrakota2023}, the package fits the GP in transformed coordinates:
\begin{equation}
x = \log m,
\qquad
z =
\begin{cases}
\log y, & y>0, \\
0, & y=0,
\end{cases}
\end{equation}
where $y$ is the observed bin count. The log transformation is used for two related
reasons. First, the prompt trident spectrum spans a wide dynamic range, and modeling
$z$ rather than $y$ turns multiplicative shape variations into approximately additive
ones. Second, $x=\log m$ makes equal distances in the GP coordinate correspond to equal
fractional changes in mass, which is the relevant scale once the detector resolution is
written in the transformed coordinate as $\sigma_x(m)\simeq \sigma_d(m)/m$. Barr and
Liu emphasize that applying logarithms to both axes improves convergence and
performance for wide, steeply falling spectra \cite{BarrLiu2025}.
The GP regression step is implemented with the
\texttt{GaussianProcessRegressor} estimator from \texttt{scikit-learn}
\cite{ScikitLearn}; the surrounding HPS analysis code supplies the histogram
handling, mass-dependent sideband masks, signal templates, profile likelihood,
and \CLs{} limit construction.

The GP is therefore used to model an unobserved smooth mean function
$z_{\mathrm{true},d}(x)$ for the log-count spectrum. In that language, ``latent'' means
only that the smooth log-rate is not observed directly; it is inferred from the
histogram counts through the transformed coordinates and the heteroscedastic variance
term. The GP mean vector is the vector of latent log-count means evaluated at the
selected training-bin coordinate. For Poisson counts $Y_i \sim
\mathrm{Pois}(\lambda_i)$, the delta method gives
$\mathrm{Var}[\log Y_i]\approx \mathrm{Var}[Y_i]/E[Y_i]^2 \approx 1/\lambda_i$, which
motivates the plug-in choice $\alpha_i\approx 1/y_i$ for occupied bins. A
heteroscedastic diagonal term $\alpha_i$ is then added to the training covariance,
\begin{equation}
\alpha_i =
\begin{cases}
1/y_i, & y_i > 0, \\
1, & y_i = 0,
\end{cases}
\label{eq:alpha-oney}
\end{equation}
\begin{equation}
K_{\mathrm{obs}} = K(X,X) + \mathrm{diag}(\bm{\alpha}),
\end{equation}
so that each $\alpha_i$ plays the role of an observation-variance term for bin $i$.
This is the quantity that the \texttt{GaussianProcessRegressor} implementation adds to
the diagonal during fitting, so in the present note it is treated as a variance model,
not as an arbitrary tuning knob.

This is a deliberate departure from the earlier Barr--Liu-inspired draft
parameterization. In the present analysis, $\alpha_i$ is treated as the latent-space
variance assignment associated with the transformed counts, while avoiding the extra
$\log y$ weighting explored in the first rough drafts. For empty bins the adopted floor
$\alpha_i=1$ is a unit variance in log-count space, not a pseudo-count of one event. It
simply prevents zero-count bins from receiving ill-defined or disproportionate weight
after the log transformation. In this high-statistics search such bins are rare, so the
floor acts mainly as a conservative legacy-stable safeguard rather than as a dominant
source of uncertainty.

\subsection{Kernel choice and CMS-motivated background modeling}
\label{sec:kernel-choice}

The baseline kernel in the current package is a single
ConstantKernel$\times$RBF model,
\begin{equation}
k(x,x') = C \exp\!\left[-\frac{(x-x')^2}{2\ell^2}\right],
\end{equation}
with broad amplitude bounds and a resolution-scaled RBF length scale. This is closely
aligned with the smooth-background prior used in contemporary CMS-oriented GPR
signal-extraction studies \cite{Gandrakota2023} and with the log-space HEP background
modeling framework emphasized in Ref.~\cite{BarrLiu2025}. The HPS implementation keeps
this kernel intentionally simple. The amplitude parameter $C$ sets the overall
covariance scale, while $\ell$ controls how rapidly the background is allowed to bend
in log-mass space. For the smooth prompt trident continuum considered here, one overall
covariance scale $C$ and one smoothness scale $\ell$ are adequate for the baseline
model: they capture the observed broad curvature without introducing additional poorly
constrained structure on signal-like mass scales.

Earlier studies preserved in the analysis history explored more complicated kernels and
piecewise procedures, but those variants are no longer the production baseline. This
note therefore documents the Constant$\times$RBF model as the authoritative kernel
choice and treats more elaborate alternatives as robustness checks rather than
competing mainline methods. The exploratory Constant$\times$Mat\'ern$(\nu=3/2)$ and
earlier piecewise-kernel studies are summarized in
Appendix~\ref{app:kernel-studies}. No production-matched alternative-kernel coverage
result is used to select the v3 baseline.

\subsection{Resolution-scaled length-scale parameterization and bounds}
\label{sec:length-scale}

Because the GP is trained in $x=\log m$, detector resolution must be translated into
log-mass units. The package defines
\begin{equation}
\sigma_x(m) = \log\!\left(m+\sigma_d(m)\right) - \log m
\simeq \frac{\sigma_d(m)}{m}
\qquad (\sigma_d \ll m),
\end{equation}
and uses this local physical scale to set the RBF bounds. In the
\path{resolution_scaled_local} policy,
\begin{equation}
\ell_{\min,d}(m) = k_{\min,d}\sigma_x(m),
\qquad
\ell_{\max,d}(m) = k_{\max,d}\sigma_x(m).
\label{eq:dataset-lengthscale-bounds}
\end{equation}
The July 2026 pre-unblinding configuration freezes the dataset-specific factors in
Table~\ref{tab:production-lengthscale-bounds}; it does not use the common
$k_{\min}=1.0$ prescription shown in earlier note drafts.

\begin{table}[H]
\centering
\begin{tabular}{lccc}
\toprule
Dataset & $k_{\min,d}$ & $k_{\max,d}$ & selection status \\
\midrule
2015 & 1.0 & 8 & frozen conservative baseline \\
2016 10\% & 0.9 & 8 & selected by functional-form closure \\
2021 1\% & 1.1 & 9 & selected calibration--reach compromise \\
\bottomrule
\end{tabular}
\caption{Resolution-scaled RBF length-scale factors in the July 2026 pre-unblinding
configuration. The factors multiply $\sigma_x(m)$ in
Eq.~\eqref{eq:dataset-lengthscale-bounds}. The labels summarize how each lower factor
is used; they do not imply that one scalar ranking metric uniquely determines all three
choices.}
\label{tab:production-lengthscale-bounds}
\end{table}

The lower factors are based on full-refit functional-form closure scans, with the
observed-spectrum result kept out of the selection criterion. For 2015, the 100-toy scan
slightly prefers $k_{\min}=1.1$ in its aggregate score, but $k_{\min}=1.0$ has similar
pull-width calibration and a weaker expected-reach proxy. The latter is therefore
retained as the conservative frozen baseline rather than described as the numerical
score winner. For 2016, the 25-toy continuation directly supports the change to
$k_{\min}=0.9$: its nonzero-injection RMS pull mean is 0.321 and its median pull width
is 0.948, compared with 0.403 and 0.889 for $k_{\min}=1.1$. It gives the smallest
aggregate calibration score among the scanned 0.9, 1.0, and 1.1 rows and reduces the
coherent apparent signal-recovery bias that motivated the follow-up. For 2021, the
corrected matched-reference scan finds small rank separations among the same three
settings. The $k_{\min}=1.1$ row has the median pull width closest to unity and the best
extraction-level $\eps^2$ proxy, while the 0.9 and 1.0 rows have slightly smaller
aggregate scores. The frozen 1.1 choice is therefore a calibration--reach tradeoff, not
a claim that it wins every diagnostic.

The upper factors remain 8 for 2015 and 2016 and 9 for 2021. They keep the optimizer
tied to a detector-scale interval while preventing effectively unbounded smoothing.
The current configurations also retain the dataset-statistics floor on
$\ell_{\max,d}$ so that the ceiling does not become artificially small where the local
resolution is exceptionally tight. Alternative upper bounds, fixed-background fits,
and the historical common-lower-bound overlays are robustness studies and are collected
in the appendices. The error-barred closure figures that directly support the frozen
choices are presented in Section~\ref{sec:toys}.

Figure~\ref{fig:gpr-window-construction} summarizes the same construction in the
language used by the implemented likelihood. The sideband bins set the log-count
training sample and its heteroscedastic bin-noise assignment, the kernel hyperparameters
are optimized only on those sidebands, and the trained GP is then conditioned to predict
the mean and covariance in the blinded mass window. That predicted count-space mean and
covariance are the background inputs passed to the local profiled signal extraction.

\begin{figure}[H]
\centering
\begin{subfigure}[t]{0.52\linewidth}
  \centering
  \graphicorplaceholder{\linewidth}{methodology_figs/gpr_training_three_step_rail.png}
  \caption{Log-count transform, noise-aware covariance training, and window prediction.
  The filled covariance boxes indicate the sideband-sideband, sideband-blind, and
  blind-blind kernel blocks used when conditioning the GP.}
\end{subfigure}
\hfill
\begin{subfigure}[t]{0.45\linewidth}
  \centering
  \graphicorplaceholder{\linewidth}{methodology_figs/blinded_mass_window_reworked.png}
  \caption{Sideband-only training and blinded-window prediction. Filled markers show
  binned count data, vertical error bars show the fixed per-bin statistical noise, and
  the horizontal reference lines are count-axis guides used to compare sideband and
  blinded-window scales.}
\end{subfigure}
\caption{Kernel and blinded-window construction used by the GPR workflow. The GP is
trained on sideband bins with fixed bin-noise variances, not on the blinded bins. The
posterior prediction in the blind window supplies the background mean and covariance
that enter the local likelihood; the shaded band in the schematic denotes the GP
posterior interval and the dashed vertical lines mark the blinded-window boundaries.}
\label{fig:gpr-window-construction}
\end{figure}

\subsection{Blind-window prediction in count space}
\label{sec:blind-window-prediction}

For each mass point, the package trains the GP only on the bins outside the training
exclusion window and predicts the posterior mean and covariance of the latent smooth
log-rate in the blinded bins. Since the fit is performed in log-count space, that
posterior is mapped back to count space through the standard lognormal moment
relations:
\begin{align}
\mu_i &= \exp\!\left(\mu^{(\log)}_i + \frac{1}{2}\Sigma^{(\log)}_{ii}\right), \\
\mathrm{Cov}(\lambda_i,\lambda_j)
&= \mu_i\mu_j\left[\exp\!\left(\Sigma^{(\log)}_{ij}\right)-1\right].
\end{align}
The resulting blind-window mean vector $\bm{b}_d(m)$ and covariance matrix
$\bm{C}_d(m)$ are the central statistical objects entering both the local extraction
fit and the \CLs\ limit. In other words, once the GP prediction has been translated
back into count space, the analysis no longer works with an abstract regression model;
it works with an explicit multivariate background prediction for the latent expected
counts in the blinded bins. The Poisson counting fluctuation of the observed
$n_{d,i}$ is handled separately by the likelihood, so $\bm{C}_d$ should be read as the
covariance of the GP-predicted background mean rather than the covariance of the
observed histogram counts themselves.

To propagate that correlated GP uncertainty into the profile likelihood, the analysis
uses the Cholesky factorization
\begin{equation}
\bm{C}_d(m)=L_d(m)L_d^\top(m),
\label{eq:Cd-cholesky}
\end{equation}
with $L_d(m)$ the lower-triangular Cholesky factor of the positive-semidefinite
blind-window covariance. Introducing an auxiliary standard-normal nuisance vector
$\bm{\theta}_d\sim\mathcal{N}(0,I)$ and the background-deformation field
$\Delta\bm{b}_d(m)=L_d(m)\bm{\theta}_d$ then gives
$\Delta\bm{b}_d\sim\mathcal{N}(0,\bm{C}_d)$. This is the standard HEP construction for
encoding correlated Gaussian-constrained nuisance variations inside a profile
likelihood model \cite{HistFactory,Cowan2011}. In the present analysis it is
especially important because the GP prediction induces nontrivial bin-to-bin
correlations across the blinded window; carrying those correlations through
$L_d\bm{\theta}_d$ lets the local significance and \CLs\ limit respond to coherent
background deformations rather than misinterpreting the GP uncertainty as independent
per-bin noise.

\subsection{Signal model and conversion to \texorpdfstring{$\eps^2$}{epsilon2}}

For each dataset $d$ and mass hypothesis $m_0$, the package first constructs a
full-range Gaussian line-shape template over the histogram binning,
\begin{equation}
w^{\mathrm{full}}_{d,i}(m_0) =
\int_{e_i}^{e_{i+1}}
\frac{1}{\sqrt{2\pi}\sigma_d(m_0)}
\exp\!\left[-\frac{(m-m_0)^2}{2\sigma_d(m_0)^2}\right]\,dm,
\end{equation}
where $e_i$ and $e_{i+1}$ are the lower and upper edges of histogram bin $i$.
The template is normalized so that
\begin{equation}
\sum_{i \in \mathrm{full\ range}} w^{\mathrm{full}}_{d,i}(m_0) = 1.
\end{equation}
The template entering the local likelihood is then the blinded-bin slice of this
full-range object,
\begin{equation}
w_{d,i}(m_0) =
\begin{cases}
w^{\mathrm{full}}_{d,i}(m_0), & i \in \mathcal{B}_d(m_0), \\
\text{not used in the local fit}, & i \notin \mathcal{B}_d(m_0).
\end{cases}
\label{eq:window-template}
\end{equation}
By the blinded-bin slice we mean exactly this restriction of the full normalized
Gaussian template to the subset of bins in $\mathcal{B}_d(m_0)$; the template is not
renormalized after the restriction.
Its blinded-bin sum is therefore
\begin{equation}
f_{\mathrm{blind},d}(m_0) = \sum_{i \in \mathcal{B}_d(m_0)} w_{d,i}(m_0) < 1,
\end{equation}
with a value near 0.90 for a $1.64\,\sigma_d$ half-width and near 0.98 for the current
$2.25\,\sigma_d$ observed-limit half-width. The signal
amplitude $A_d$ is defined as the \emph{total} local Gaussian signal yield associated
with the mass hypothesis, not merely the yield inside the blind window. The expected
signal contribution entering bin $i$ of the local likelihood is therefore
\begin{equation}
s_{d,i}(A_d;m_0) = A_d\,w_{d,i}(m_0),
\label{eq:signal-bin-contribution}
\end{equation}
so the total expected signal actually visible inside the blind window is
$f_{\mathrm{blind},d}(m_0)\,A_d$.

This convention is more physical than blind-window renormalization because it preserves
the fact that a finite blind region never contains the entire Gaussian line shape. It
also keeps the signal parameter used in the fit aligned with the full-histogram
injection model used in the procedural toy studies.

The conversion between signal yield and dark-photon coupling follows the standard
fixed-target relation between $A^\prime$ production and radiative trident production
\cite{EssigEtAl2011,HPS2015DarkPhoton,HPS2016PromptLong}. In the present notation,
\begin{equation}
\frac{dN_{A^\prime}}{dm}\bigg|_{m=m_{A^\prime}}
= \frac{3\pi}{2N_{\mathrm{eff}}\alpha_{\mathrm{EM}}}
\; m_{A^\prime}\,\eps^2\,
\frac{dN_{\gamma^\ast}}{dm}\bigg|_{m=m_{A^\prime}},
\end{equation}
where $\alpha_{\mathrm{EM}}$ is the fine-structure constant, the factor $3\pi/2$ is the
standard narrow-width coefficient relating on-shell $A^\prime$ production to radiative
trident production, and $N_{\mathrm{eff}}$ counts the number of kinematically open
visible decay channels contributing to the total $A^\prime$ width. Below the dimuon
threshold, only the $e^+e^-$ channel is open and the minimal-model conversion has
$N_{\mathrm{eff}}=1$. The present GPR implementation keeps this $e^+e^-$ normalization
as the analysis convention for the prompt electron-channel yield; minimal dark-photon
interpretations above $2m_\mu$ require updating the visible branching factor to include
the newly open $\mu^+\mu^-$ channel. Writing the
radiative differential rate as
\begin{equation}
\frac{dN_{\gamma^\ast}}{dm}
= f_{\mathrm{rad},d}(m)\,\rho_d(m),
\end{equation}
where $\rho_d(m)$ is the differential prompt background rate in the selected
invariant-mass spectrum, the package converts between $A_d$ and $\eps^2$ through
\begin{equation}
\eps^2 =
\frac{2\alpha_{\mathrm{EM}}\,A_d}
{3\pi\,m\,f_{\mathrm{rad},d}(m)\,\rho_d(m)},
\label{eq:eps2-conversion}
\end{equation}
where $\rho_d(m)$ is estimated from the observed histogram as a local counts-per-GeV
density in the same mass interval used for the nominal blind window,
\begin{equation}
\rho_d(m)\equiv
\frac{N_d\!\left[m-n_{\mathrm{blind}}\sigma_d(m),\,m+n_{\mathrm{blind}}\sigma_d(m)\right]}
{2\,n_{\mathrm{blind}}\sigma_d(m)},
\end{equation}
with $n_{\mathrm{blind}}=1.64$ retained as the current prompt-density normalization
window even when the extraction blind window is widened to $2.25\,\sigma_d$. This keeps
the yield-to-$\eps^2$ conversion tied to the historically used local density scale while
the likelihood itself can recover the wider signal tails. The code exposes
\texttt{eps2\_density\_nsigma} as an audit knob, so alternative density-window
conventions can be reproduced when needed. It is convenient to define the
dataset-specific conversion factor
\begin{equation}
K_d(m) \equiv A_d(\eps^2 = 1)
= \frac{3\pi\,m\,f_{\mathrm{rad},d}(m)\,\rho_d(m)}{2\alpha_{\mathrm{EM}}},
\label{eq:Kd-conversion}
\end{equation}
so that a shared coupling $\eps^2$ predicts a dataset-specific total signal yield
$A_d = K_d(m)\eps^2$.

Unless stated otherwise, the methodology section uses the nominal radiative fraction
$f_{\mathrm{rad},d}(m)$. The optional penalty that propagates normalization uncertainty
into $\eps^2$ is treated as a systematic variation rather than as part of the core GP
fit and is discussed separately in Section~\ref{sec:radfrac-penalty}. The numerical
penalty values are therefore result-configuration choices, not ingredients of the
background interpolation or local extraction fit. When that systematic option is
enabled for dataset $d$, the effective radiative fraction is
\begin{equation}
f_{\mathrm{rad},d}^{\mathrm{eff}}(m)
= \left(1-\delta_{f,d}\right) f_{\mathrm{rad},d}(m),
\end{equation}
so the same signal-yield limit maps to a weaker $\eps^2$ limit.

\subsection{Statistical inference in a single dataset}

For a fixed dataset $d$ and mass hypothesis $m$, let $\bm{n}_d$ denote the vector of
observed blinded-bin counts, $\bm{b}_d$ the GP-predicted blinded-bin mean,
$\bm{C}_d=L_dL_d^\top$ the corresponding covariance of
Section~\ref{sec:blind-window-prediction} and Eq.~\eqref{eq:Cd-cholesky}, and
$\bm{w}_d(m)$ the
non-renormalized blinded-bin slice of the full Gaussian template defined in
Eq.~\eqref{eq:window-template}. The profiled local likelihood is
\begin{equation}
\mathcal{L}_d(A_d,\bm{\theta}_d)
=
\left[
\prod_{i=1}^{N_d(m)}
\mathrm{Pois}\!\left(
n_{d,i}\,\middle|\,
b_{d,i} + (L_d\bm{\theta}_d)_i + A_d w_{d,i}(m)
\right)
\right]
\exp\!\left(-\frac{1}{2}\bm{\theta}_d^\top\bm{\theta}_d\right),
\label{eq:singlelike}
\end{equation}
where the Gaussian nuisance vector $\bm{\theta}_d\sim\mathcal{N}(0,I)$ induces the
correlated background-deformation field $L_d\bm{\theta}_d$ defined below
Eq.~\eqref{eq:Cd-cholesky}, and $N_d(m)$ is the number of bins in the local blinded
set $\mathcal{B}_d(m)$. The signal contribution entering the Poisson mean is exactly
the blinded-bin restriction of Eq.~\eqref{eq:signal-bin-contribution}.

At fixed $m$, profiling Eq.~\eqref{eq:singlelike} with no physical sign constraint on
$A_d$ yields the reviewer-facing extraction estimator $\hat A_d^{\,\mathrm{unc}}$ and
its local uncertainty $\sigma_{A,d}$. This unconstrained fit is used for closure,
linearity, and pull studies because it allows downward background fluctuations to be
captured by the estimator itself rather than being clipped at zero. For discovery and
production upper limits, however, the parameter of interest is restricted to the
physical region $A_d\ge0$, and the corresponding physical best fit is
\begin{equation}
\hat A_d = \max\!\left(0,\hat A_d^{\,\mathrm{unc}}\right).
\end{equation}
We denote by $\hat{\bm{\theta}}_d$ the nuisance-parameter maximizer at the
unconditional best fit and by $\hat{\hat{\bm{\theta}}}_d(A_d)$ the conditional maximum
likelihood estimator of $\bm{\theta}_d$ for fixed $A_d$.

\subsubsection{Local significance}

The single-dataset discovery statistic is the bounded profile-likelihood ratio
\begin{equation}
q_{0,d}(m) = -2\ln
\frac{\mathcal{L}_d\!\left(A_d=0,\hat{\hat{\bm{\theta}}}_d(0)\right)}
{\mathcal{L}_d\!\left(\hat A_d,\hat{\bm{\theta}}_d\right)},
\label{eq:q0-single}
\end{equation}
with $\hat A_d\ge0$. The local one-sided discovery $p$-value and the corresponding
Gaussian significance are then obtained with the asymptotic Cowan-et-al.\ mapping
\cite{Cowan2011}. The resulting local-significance definition is explicitly
one-sided and discovery facing: downward fluctuations are allowed in the unconstrained
fit used for closure studies, but they do not generate a positive local discovery
significance. This profiled multi-bin test replaces the older sum-count approximation
used in the earliest drafts and is the only local-significance definition treated as
current in this note.

\subsubsection{\texorpdfstring{\CLs}{CLs} upper limits}

To test a fixed non-negative signal amplitude $A_d$, the production limit setting uses
the bounded one-sided profile-likelihood statistic
\begin{equation}
\tilde q_{A_d}(m)=
\begin{cases}
0, &
\hat A_d^{\,\mathrm{unc}} > A_d, \\
-2\ln
\dfrac{\mathcal{L}_d\!\left(A_d,\hat{\hat{\bm{\theta}}}_d(A_d)\right)}
{\mathcal{L}_d\!\left(\hat A_d^{\,\mathrm{unc}},\hat{\bm{\theta}}_d\right)}, &
0\le \hat A_d^{\,\mathrm{unc}}\le A_d, \\
-2\ln
\dfrac{\mathcal{L}_d\!\left(A_d,\hat{\hat{\bm{\theta}}}_d(A_d)\right)}
{\mathcal{L}_d\!\left(0,\hat{\hat{\bm{\theta}}}_d(0)\right)}, &
\hat A_d^{\,\mathrm{unc}} < 0,
\end{cases}
\label{eq:qmu-single}
\end{equation}
which is the standard bounded Cowan construction for upper limits \cite{Cowan2011}.
The corresponding modified-frequentist tail areas are
\begin{equation}
\mathrm{CL}_{s+b}(A_d;m)
=
P\!\left(
\tilde q_{A_d}\ge \tilde q_{A_d}^{\mathrm{obs}}
\middle|\,
A_d+\mathrm{bkg}
\right),
\qquad
\mathrm{CL}_{b}(A_d;m)
=
P\!\left(
\tilde q_{A_d}\ge \tilde q_{A_d}^{\mathrm{obs}}
\middle|\,
\mathrm{bkg}
\right),
\end{equation}
and
\begin{equation}
\mathrm{CL}_{s}(A_d;m)=
\frac{\mathrm{CL}_{s+b}(A_d;m)}{\mathrm{CL}_{b}(A_d;m)}.
\label{eq:cls-definition}
\end{equation}
Dividing by $\mathrm{CL}_b$ prevents a downward background fluctuation from producing
an artificially aggressive exclusion in a region with little intrinsic sensitivity
\cite{Junk1999,Read2002}. At confidence level $1-\alpha$, the single-dataset upper
limit is
\begin{equation}
A_{\mathrm{UL},d}(m;\alpha)
=\min\{A_d:\mathrm{CL}_s(A_d;m)\le\alpha\}.
\label{eq:aul-single}
\end{equation}
The v3 candidate uses $\alpha=0.10$, corresponding to a 90\% CL upper limit; retained
historical 95\% studies use $\alpha=0.05$. The package reports the limit both in
signal-yield space and, after applying the dataset-specific conversion factor of
Eq.~\eqref{eq:Kd-conversion}, in $\eps^2$ space.

\subsection{Asymptotic and toy calibration}

The local-significance mapping introduced above is treated asymptotically throughout
the production scan. For upper limits, the same bounded statistic
$\tilde q_{A_d}$ is calibrated in one of two ways. In asymptotic mode, the
background-only Asimov data set is used to evaluate the reference quantity
$\tilde q_{A_d,A}$, so that
\begin{equation}
\mathrm{CL}_{s+b}(A_d;m)
=
1-\Phi\!\left(\sqrt{\tilde q_{A_d}^{\mathrm{obs}}}\right),
\qquad
\mathrm{CL}_{b}(A_d;m)
=
\Phi\!\left(
\sqrt{\tilde q_{A_d,A}}-\sqrt{\tilde q_{A_d}^{\mathrm{obs}}}
\right).
\label{eq:cls-asymptotic}
\end{equation}
In toy mode, the same statistic $\tilde q_{A_d}$ is calibrated directly with
background-only and signal-plus-background pseudoexperiments rather than by switching
to a different ordering variable. The shared-coupling combination described below uses
the identical logic after replacing the single-dataset amplitude $A_d$ with the common
parameter of interest $\eps^2$.

The pseudoexperiment program remains central even with the asymptotic formulae in
place because the HPS problem is not a textbook one-bin counting limit. The local
likelihood combines a multi-bin Poisson term, a GP-conditioned blind-window
covariance, a non-renormalized Gaussian signal template, and a mass-dependent
yield-to-$\eps^2$ conversion. The pull, coverage, and $Z$-calibration studies of
Section~\ref{sec:toys} therefore validate not only code closure but also the domain in
which the asymptotic summaries can be interpreted with confidence.

\subsection{Observed-limit toy diagnostics}

The analysis uses two distinct families of $p$-values, and the distinction is
important for internal review. The first family is the discovery statistic of
Eq.~\eqref{eq:q0-single}, summarized through $q_0(m)$, $p_0(m)$, and
$Z_{\mathrm{local}}(m)$. The second family is a set of reviewer-facing diagnostics
derived from the background-only \CLs{} limit ensemble at each mass. For
$N_{\mathrm{toys}}$ background-only pseudoexperiments and an observed limit
$\eps^2_{\mathrm{UL},\mathrm{obs}}(m)$ at the confidence level under study, those
quantities are
\begin{align}
p_{\mathrm{strong}}(m)
&=
\frac{1}{N_{\mathrm{toys}}}
\sum_{t=1}^{N_{\mathrm{toys}}}
\mathbf{1}\!\left[
\eps^2_{\mathrm{UL},t}(m)\le \eps^2_{\mathrm{UL},\mathrm{obs}}(m)
\right], \\
p_{\mathrm{weak}}(m)
&=
\frac{1}{N_{\mathrm{toys}}}
\sum_{t=1}^{N_{\mathrm{toys}}}
\mathbf{1}\!\left[
\eps^2_{\mathrm{UL},t}(m)\ge \eps^2_{\mathrm{UL},\mathrm{obs}}(m)
\right], \\
p_{\mathrm{two}}(m)
&=
\min\!\left[1,\,
2\min\!\left(p_{\mathrm{strong}}(m),\,p_{\mathrm{weak}}(m)\right)\right].
\end{align}
Here $p_{\mathrm{strong}}$ quantifies how unusually \emph{strong} the observed limit
is relative to the background-only toy ensemble, while $p_{\mathrm{weak}}$ quantifies
how unusually \emph{weak} it is. The two-sided quantity $p_{\mathrm{two}}$ is retained
only as a symmetric reviewer diagnostic; because it responds to both unusually strong
and unusually weak limits, it is less targeted than the one-sided discovery or
exclusion tests and is not used as a search statistic. Figure~\ref{fig:pvalue-tail-schematic}
gives the compact fixed-mass definition of these quantities, while
Figure~\ref{fig:pvalue-tail-examples} shows representative strong-limit, weak-limit,
and bounded median-case examples from the same toy-limit logic.

\begin{figure}[H]
\centering
\graphicorplaceholder{0.96\linewidth}{methodology_figs/pvalue_tail_schematic.png}
\caption{Illustrative toy-limit diagnostic at fixed mass. The left panel shows a
background-only ensemble of observed \CLs{} upper limits, and the right panel shows the
corresponding empirical cumulative distribution. The vertical line marks the observed
limit. The lower-tail area defines $p_{\mathrm{strong}}$, the upper-tail area defines
$p_{\mathrm{weak}}$, and
$p_{\mathrm{two}}=\min[1,2\min(p_{\mathrm{strong}},p_{\mathrm{weak}})]$ is used only
as a symmetric diagnostic summary.}
\label{fig:pvalue-tail-schematic}
\end{figure}

\begin{figure}[H]
\centering
\graphicorplaceholder{0.98\linewidth}{methodology_figs/pvalue_tail_examples.png}
\caption{Representative fixed-mass observed-limit toy diagnostics. The left panel
shows an unusually strong observed limit relative to the background-only toy
ensemble, the middle panel shows an unusually weak observed limit, and the right
panel shows a near-median case where the bounded symmetric summary saturates at
$p_{\mathrm{two}}=1$. These are reviewer-facing limit diagnostics, not discovery
statistics.}
\label{fig:pvalue-tail-examples}
\end{figure}

These quantities are therefore limit diagnostics rather than discovery-significance
measures in the sense of $p_0$. If the observed limit lies near the median of the
background-only ensemble, then roughly half of the toys lie above it and half lie
below it, making both $p_{\mathrm{strong}}$ and $p_{\mathrm{weak}}$ naturally
$\mathcal{O}(0.5)$. With $N_{\mathrm{toys}}=\num{10000}$, the smallest non-zero
empirical tail probability is still only of order $10^{-4}$, so these toy-limit
diagnostics cannot by themselves establish or refute a high-significance local excess
in the discovery statistic defined by the one-sided local $p_0$ mapping above.

\subsection{Shared-\texorpdfstring{$\eps^2$}{epsilon2} combination across datasets}
\label{sec:combination}

At a common mass hypothesis $m$, each active dataset contributes a distinct blinded-bin
count vector, GP background prediction, covariance matrix, and Gaussian signal
template. Let the set of datasets that cover the hypothesis be $\mathcal{D}(m)$. For
each $d \in \mathcal{D}(m)$, define
\begin{equation}
\bm{n}_d(m),\qquad
\bm{b}_d(m),\qquad
\bm{C}_d(m)=L_dL_d^\top,\qquad
\bm{w}_d(m),\qquad
K_d(m),
\end{equation}
where $\bm{w}_d(m)$ is the blinded-bin template slice of
Eq.~\eqref{eq:window-template}, $K_d(m)$ is the yield-per-unit-$\eps^2$ conversion of
Eq.~\eqref{eq:Kd-conversion}, and $L_d$ is the Cholesky factor of the GP covariance in
that dataset. When the radiative-fraction penalty of Eq.~\eqref{eq:frad-penalty} is
enabled, the conversion becomes
\begin{equation}
K_d^{\mathrm{eff}}(m) = \left(1-\delta_{f,d}\right)K_d(m),
\label{eq:kdeff}
\end{equation}
so the expected signal vector per unit $\eps^2$ is reduced coherently in every active
dataset. The per-dataset likelihood is then
\begin{equation}
\mathcal{L}_d(\eps^2,\bm{\theta}_d)
=
\left[
\prod_{i=1}^{N_d(m)}
\mathrm{Pois}\!\left(
n_{d,i}\,\middle|\,
b_{d,i} + (L_d\bm{\theta}_d)_i + \eps^2 K_d^{\mathrm{eff}}(m) w_{d,i}(m)
\right)
\right]
\exp\!\left(-\frac{1}{2}\bm{\theta}_d^\top\bm{\theta}_d\right).
\label{eq:dataset-poi-like}
\end{equation}
The baseline combined likelihood is the product over datasets,
\begin{equation}
\mathcal{L}_{\mathrm{comb}}(\eps^2,\{\bm{\theta}_d\})
= \prod_{d \in \mathcal{D}(m)} \mathcal{L}_d(\eps^2,\bm{\theta}_d).
\label{eq:combinedlike-product}
\end{equation}
Equivalently, one may concatenate all blinded-bin vectors into a single object. In the
following equations, $\mathbin{\Vert}$ denotes vector concatenation, e.g.
$\bm{n}_1\mathbin{\Vert}\bm{n}_2$:
\begin{align}
\bm{n} &=
\bm{n}_1\mathbin{\Vert}\bm{n}_2\mathbin{\Vert}\cdots, \\
\bm{b} &=
\bm{b}_1\mathbin{\Vert}\bm{b}_2\mathbin{\Vert}\cdots, \\
\bm{s}_{\mathrm{unit}} &=
\left(K_1^{\mathrm{eff}}\bm{w}_1\right)
\mathbin{\Vert}
\left(K_2^{\mathrm{eff}}\bm{w}_2\right)
\mathbin{\Vert}\cdots, \\
\bm{\theta} &=
\bm{\theta}_1\mathbin{\Vert}\bm{\theta}_2\mathbin{\Vert}\cdots, \\
\bm{C}_{\mathrm{comb}} &=
\mathrm{blockdiag}(\bm{C}_1,\bm{C}_2,\ldots),
\qquad
L_{\mathrm{comb}}=\mathrm{blockdiag}(L_1,L_2,\ldots),
\end{align}
so that the total expected mean is
\begin{equation}
\bm{\lambda}(\eps^2,\bm{\theta})
= \bm{b} + L_{\mathrm{comb}}\bm{\theta} + \eps^2\bm{s}_{\mathrm{unit}},
\qquad
\bm{C}_{\mathrm{comb}} = L_{\mathrm{comb}}L_{\mathrm{comb}}^\top.
\end{equation}
The bounded upper-limit statistic for the shared coupling,
$\tilde q_{\eps^2}(m)$, is defined exactly as in Eq.~\eqref{eq:qmu-single} after
replacing the single-dataset amplitude with the common parameter of interest
$\eps^2$ and the single-dataset likelihood with
$\mathcal{L}_{\mathrm{comb}}$. Two points are essential. First, the baseline
combination shares only the parameter of interest $\eps^2$. The nuisance vectors
$\bm{\theta}_d$ are independent between datasets, and no common radiative-fraction,
luminosity, or resolution nuisance parameters are profiled across runs in the current
production model. Second, the combination is performed at the likelihood level, not by
averaging precomputed upper-limit curves. At fixed $\eps^2$,
\begin{equation}
-2\ln\lambda_{\mathrm{comb}}(\eps^2)
=
\sum_{d \in \mathcal{D}(m)} -2\ln\lambda_d(\eps^2),
\end{equation}
after profiling the independent nuisance blocks. The sensitivity gain therefore comes
from the additive information content of statistically independent channels, each with
its own detector response and normalization factor, rather than from any post hoc
merging of already-compressed outputs. This is the standard common-POI structure used
in multi-channel HEP combinations \cite{HistFactory,ATLASCMSHiggs2016}. The combined
upper limit is set directly on the shared coupling,
\begin{equation}
\eps^2_{\mathrm{UL}}(m;\alpha)
= \min\{\eps^2 : \mathrm{CL}_s(\eps^2)\le\alpha\},
\end{equation}
with $\alpha=0.10$ for the v3 90\% CL candidate. The combined local significance is
evaluated from the same profiled likelihood
with the common signal vector $\eps^2 \bm{s}_{\mathrm{unit}}$. By contrast, the
inverse-variance expressions shown later in the search-power figures are auxiliary
interpretations of combined sensitivity, not the defining inference.

\subsection{Global significance and the look-elsewhere treatment}

The physics-facing local significance is the minimum one-sided $p_0$ in the scan,
\begin{equation}
p_{0,\min}^{\mathrm{obs}} = \min_m p_0^{\mathrm{obs}}(m),
\qquad
q_{0,\max}^{\mathrm{obs}} = \max_m q_0^{\mathrm{obs}}(m).
\end{equation}
The cleanest global-significance definition is then scan based:
\begin{equation}
p_{\mathrm{global}}^{\mathrm{toy}}
=
\frac{1}{N_{\mathrm{toys}}}
\sum_{t=1}^{N_{\mathrm{toys}}}
\mathbf{1}\!\left[
q_{0,\max}^{(t)} \ge q_{0,\max}^{\mathrm{obs}}
\right],
\label{eq:pglobal-toy}
\end{equation}
with
\begin{equation}
Z_{\mathrm{global}}^{\mathrm{toy}} = \Phi^{-1}\!\left(1-p_{\mathrm{global}}^{\mathrm{toy}}\right).
\end{equation}
When a dense full-scan toy calibration is not yet frozen, the workflow monitors the
look-elsewhere effect through an effective-trials approximation based on the scan span
and detector resolution \cite{GrossVitells2010}. In that representation,
\begin{equation}
p_{\mathrm{global}}^{\mathrm{Sidak}}
=
1-\left(1-p_{0,\min}^{\mathrm{obs}}\right)^{N_{\mathrm{eff}}},
\qquad
Z_{\mathrm{global}}^{\mathrm{Sidak}}
=
\Phi^{-1}\!\left(1-p_{\mathrm{global}}^{\mathrm{Sidak}}\right).
\label{eq:sidak-global}
\end{equation}
For small $p_{0,\min}^{\mathrm{obs}}$, Eq.~\eqref{eq:sidak-global} reduces to the
familiar $p_{\mathrm{global}}\simeq N_{\mathrm{eff}}p_{0,\min}$ form.

For the present studies, Eq.~\eqref{eq:sidak-global} provides the global, or
look-elsewhere-corrected, $p$-value quoted with each local scan. The effective number of
trials is evaluated from the stated scan range and the mass-dependent resolution, so
the approximation accounts for the strong correlation between adjacent
\SI{1}{MeV}-spaced hypotheses rather than treating every scan point as independent.
This is sufficient for interpreting the current pre-unblinding scans and is shown
alongside the local $p$-values in Section~\ref{sec:results}. A background-only ensemble
of the frozen search maximum, $q_{0,\max}$, remains the preferred calibration for any
future discovery-level claim; when that ensemble is available,
Eq.~\eqref{eq:pglobal-toy} supersedes the analytic approximation. The distinction is
therefore one of calibration precision, not an absence of a global interpretation in
v3.

\subsection{Signal absorption and alternative background-profile treatments}

The production baseline documented here is the additive Gaussian-profile treatment of
blind-window background uncertainty. That choice is natural because the GP delivers an
additive mean vector and covariance matrix in the blinded bins, and the entire current
toy-validation program is built around that object. At the same time, the remaining
signal-absorption question should be stated explicitly. The newly adopted
leakage-retaining signal model removes one obvious source of bias associated with
blind-window renormalization, but it does not by itself prove that the full procedural
toy workflow is free of absorption once the GP is refit toy by toy.

The most likely residual sources are the refitting freedom of the GP hyperparameters,
especially the effective length scale, and the flexibility of the Gaussian-profile
background nuisance itself. For that reason, the current note keeps the
Gaussian-profile model as the mainline baseline while reserving fixed-background and
alternative Poisson-style treatments as targeted robustness studies. Those alternatives
are not presented as equally mature production choices. They become necessary
follow-up tests if production-faithful pseudoexperiments demonstrate that the current
nuisance model absorbs an unacceptable fraction of injected signal.

\subsection{Dimuon-threshold interpretation}
\label{sec:methodology-dimuon-threshold}

The yield-to-coupling conversion in Eq.~\eqref{eq:eps2-conversion} is an
electron-channel normalization: the fitted resonance amplitude is the
$A^\prime\to e^+e^-$ yield visible in the prompt invariant-mass spectrum. For the
minimal kinetically mixed dark photon, this is also the full visible branching
interpretation below the dimuon threshold, $m_{A^\prime}<2m_\mu=\SI{211.3167}{MeV}$,
where the $\mu^+\mu^-$ final state is kinematically closed. This is the same
fixed-target convention used in the published HPS prompt searches
\cite{EssigEtAl2011,HPS2015DarkPhoton,HPS2016PromptLong}. Above threshold, the same
electron-channel yield corresponds to a weaker minimal-model limit on $\eps^2$ because
the total width includes the newly open $A^\prime\to\mu^+\mu^-$ mode
\cite{Holdom1986,Fabbrichesi2020DarkPhoton}.

Neglecting the electron mass and remaining below the two-pion hadronic threshold, the
charged-lepton width ratio is
\begin{equation}
R_{\mu/e}(m_{A^\prime})
\equiv
\frac{\Gamma(A^\prime\to\mu^+\mu^-)}
{\Gamma(A^\prime\to e^+e^-)}
=
\sqrt{1-\frac{4m_\mu^2}{m_{A^\prime}^2}}\,
\left(1+\frac{2m_\mu^2}{m_{A^\prime}^2}\right),
\qquad m_{A^\prime}>2m_\mu,
\label{eq:dimuon-width-ratio}
\end{equation}
with $R_{\mu/e}=0$ below threshold. The inverse electron branching fraction used for
the minimal-model reinterpretation is therefore
\begin{equation}
N_{\mathrm{eff}}^{\mathrm{BR}}(m_{A^\prime})
\equiv
\frac{1}{\mathcal{B}(A^\prime\to e^+e^-)}
=
1+R_{\mu/e}(m_{A^\prime}),
\label{eq:dimuon-neff}
\end{equation}
and an electron-channel coupling limit is mapped to the minimal visible-dark-photon
convention by
\begin{equation}
\eps^2_{\mathrm{min}}(m_{A^\prime})
=
N_{\mathrm{eff}}^{\mathrm{BR}}(m_{A^\prime})\,\eps^2_{ee}(m_{A^\prime}).
\label{eq:dimuon-eps2-map}
\end{equation}
This operation does not change the local fit, the signal-yield upper limit, the
\CLs{} value, or the local $p_0$ scan; it changes only the coupling interpretation of
an electron-channel result. Figure~\ref{fig:methodology-dimuon-factor} shows the
factor applied in the v3 90\% CL candidate obtained with \CLs{}. The factor is unity below
$2m_\mu$, rises continuously above threshold, and reaches
$N_{\mathrm{eff}}^{\mathrm{BR}}=1.725$ at \SI{250}{MeV}
($\mathcal{B}_{ee}=0.580$).

The scan range used for a look-elsewhere statement must match the range of the quoted
physics claim \cite{GrossVitells2010}. If the minimal-dark-photon exclusion is quoted
only below \SI{210}{MeV}, the final scan-level toy or effective-trials calibration
should use that same below-threshold range. If the high-mass electron-channel scan is
promoted to a \SIrange{170}{250}{MeV} minimal-model result, the corresponding
look-elsewhere calibration should include the full extended range as well.

\begin{figure}[H]
\centering
\graphicorplaceholder{0.78\linewidth}{final_limit_projection_figs/90cls/dimuon/dimuon_factor_vs_mass.png}
\caption{Minimal-model branching factor used to reinterpret electron-channel
$\eps^2$ limits above the dimuon threshold. The red curve shows
$N_{\mathrm{eff}}^{\mathrm{BR}}=1/\mathcal{B}(A^\prime\to e^+e^-)$ and the blue dashed
curve shows $\mathcal{B}(A^\prime\to e^+e^-)$.}
\label{fig:methodology-dimuon-factor}
\end{figure}
